\section{The conditional theorem}
\label{sec:landing}

We can now state the main result.  \Cref{fig:landing-chain-diagram}
shows its logical structure: which hypotheses enter the proof of the
scalar identity, and which are carried along.

\begin{figure}[tbp]
\centering
\begin{tikzpicture}[
  font=\small,
  input/.style={draw, rounded corners=2pt, align=center, minimum width=4.7cm,
    minimum height=1.35cm, fill=black!3},
  box/.style={draw, rounded corners=2pt, align=center, minimum width=3.0cm,
    minimum height=1.0cm, fill=black!3},
  result/.style={draw, rounded corners=2pt, align=center, minimum width=4.0cm,
    minimum height=1.35cm, fill=black!8},
  arrow/.style={-{Latex[length=2mm]}, thick}
]
\node[input] (inputs) at (-4.9,0.25) {shell-closure record\\
  three recognition records\\headline imports};
\node[box] (master) at (0,0.25) {closure master\\theorem applies};
\node[result] (landing) at (4.9,0.25) {\(\rhoE=0\)\\
  {\scriptsize readout equivalence}\\Strong-BSD identity};

\draw[arrow] (inputs) -- (master);
\draw[arrow] (master) -- (landing);
\node[align=center, font=\footnotesize] at (0,-1.15)
  {conditional on the supplied records, the shell residual vanishes and the scalar identity follows};
\end{tikzpicture}
\caption{The conditional landing chain from the shell-closure record,
recognition records, and headline imports to the Strong-BSD scalar identity.}
\label{fig:landing-chain-diagram}
\end{figure}


\begin{theorem}[Conditional Strong BSD]\label{thm:closure:strong-bsd-conditional}
Let $\SelBSD$ be a shell with curve $E$ and analytic rank $r$, and let
$C$ be its closure operator.  Assume:
\begin{enumerate}
\item the recognition predicate $\PiBSD(E)$
  (\Cref{def:closure:pi-bsd});
\item the Cassels--Tate input $\chiCTp$ at $E$ and some prime $p$, and the
  Selmer-complex pairing input $\AofE$ at $E$
  (\Cref{def:closure:chi-ct-p-import,def:closure:a-e-import});
\item if $r\ge2$, the archimedean comparison input at $E$, of rank $r$
  (\Cref{def:closure:beilinson-import}).
\end{enumerate}
Then
\[
  \frac{L^{(r)}(E,1)}{r!}
  =
  \frac{\#\Sha(E/\Q)\,\Reg_{\mathrm{NT}}(E)\,\Omega_E\,\Tam(E)}
       {\#E(\Q)_{\tors}^{\,2}},
\]
equivalently $\rhoE=0$, and the asserted components of the inputs in
(2) and (3) hold.%
\footnote{The formalization records this theorem as a conditional
schema, because it returns the components of the inputs in (2) and (3)
by projection.  Its scalar conclusion, however, is derived: Lean proves
it from items (3) and (4) of $\PiBSD(E)$ by the algebra in the proof
below.  A compiled dependency check confirms that the two declarations
isolating this algebra depend on no other declaration of the project,
and that neither the derivation from $\PiBSD(E)$ nor the full theorem
depends on \Cref{thm:closure:chi-ct-p-comparison},
\Cref{thm:closure:oc-eta-formula}, \Cref{thm:closure:composite-signature},
the statements of \Cref{sec:aor-instance}, or the supplied applicability,
computability, falsifiability and ablation fields of the shell.}
\end{theorem}

\begin{proof}
If $r\in\{0,1\}$, item (4) of $\PiBSD(E)$ is the identity $L^*=R_E$.
If $r\ge2$, item (3) supplies a $\GammaGZ$ record whose factors are those
of the shell and whose constant is $\kapr=\#\Sha/d$, with $d$ the squared
torsion order.  Its hypothesis $\psi_-(E,r)=0$ says
$L^*-\kapr\Reg\,\Omega\,\Tam=0$, hence $L^*=\kapr\Reg\,\Omega\,\Tam$ by the
first law of the shell.  Substituting $\kapr=\#\Sha/d$ and applying the
second law $(a/d)\cdot b=(a\cdot b)/d$ three times gives
\[
  L^*
  =\Bigl(\frac{\#\Sha}{d}\Bigr)\Reg\,\Omega\,\Tam
  =\frac{\#\Sha\,\Reg\,\Omega\,\Tam}{d}
  =R_E .
\]
In both cases $\rhoE=0$ by \Cref{thm:closure:pi-bsd-iff-strong-bsd}.  The
remaining assertions are components of the inputs.
\end{proof}

Three comments make the content of the theorem precise.

\emph{What the proof uses.}  The scalar conclusion is derived from the
rank-split part of $\PiBSD(E)$ alone.  The local hypotheses $\GammaPD$
and $\GammaSP$ and the comparison inputs are hypotheses of the theorem,
because they belong to the arithmetic setting in which the recognition
predicate is meant, but the scalar identity does not depend on them.

\emph{How strong the hypothesis is.}  The rank-split scalar premise
contains the identity in rank $0$ and $1$, and expresses it through
normalized fixity in rank at least $2$: $\GammaGZ$ with the normalization
$\kapr=\#\Sha/\#E(\Q)_{\tors}^{\,2}$ is the identity rewritten.  The
remaining hypotheses supply additional records, which the identity does
not reconstruct.  The theorem is therefore not an independent derivation of
Strong BSD from weaker arithmetic statements, and it is not a statement
about minimal hypotheses.  Its content is the exact translation: the
identity is equivalent to the vanishing of $\rhoE$
(\Cref{thm:closure:pi-bsd-iff-strong-bsd}) and, when the scalars form a
field with $2\neq0$, to the fixedness of $x_E$ under the averaging closure
(\Cref{thm:closure:scalar-closure-translation}),
and it follows from the recognition predicate by explicit algebra, with
every arithmetic input named and scoped.

\emph{That the hypotheses are consistent and not vacuous.}  The theorem
has been instantiated on shells with scalars in $\Q$, a label in place of
a curve, all factors equal to $1$, and rank $0$ or $2$.  With leading
coefficient $1$, all hypotheses hold and the conclusion follows; in rank
$0$ no archimedean input is required, and none can be supplied.  With
leading coefficient $2$ and everything else unchanged, the recognition
predicate fails, since it would force $2=1$.  These are formal checks of
the statement, not examples of elliptic curves.

\begin{nonclaim}\label{nonclaim:closure:landing-conditional}
This is not a proof of the Birch--Swinnerton-Dyer conjecture, in its weak
or strong form, for any curve.  The theorem derives the Strong BSD
identity for $E$ from hypotheses whose scalar part contains it in rank at
most one and expresses it through normalized fixity in rank at least two.
We know of no theorem supplying the leading-term identity in rank at
least two, or the main conjectures at additive primes $2$ and $3$
(\Cref{rmk:closure:t-e12-gap,rmk:closure:t-bad-gap}).
\end{nonclaim}
